%
% Tex input file for CGNS Overview document
%
% To generate a DVI file, then a PostScript file named overview.ps:
%
%     latex overview.tex
%     dvips overview.dvi -o
%
% To generate a PDF file named overview.pdf:
%
%     pdflatex overview.tex


\documentclass[twoside,fleqn]{article}

%
% Packages
%
%\usepackage{times}			% use PostScript Times-Roman for text
\usepackage{indentfirst}		% indent first paragraph in sections
%\usepackage[first,light]{draftcopy}	% add DRAFT watermark
\usepackage[tbtags]{amsmath}		% use ams math package
\usepackage{graphicx}			% for including eps files
%\usepackage{subfigure}			% for sub-figures
\usepackage[bf]{caption}		% for more flexible captions
   \setlength{\captionmargin}{18pt}	% caption margins
\usepackage{array}			% use new array and table package
   \setlength{\extrarowheight}{2pt}	% increase row heights in tables
\usepackage{longtable}			% use longtable package
\usepackage{alltt}			% verbatim input with font changes
%\usepackage{lscape}			% for landscape pages
%\usepackage{mdwtab}			% use mdwtab array and table package
%\usepackage{tabls}			% access \hline[dimen] feature
%\usepackage{dcolumn}			% allow table alignment at dec points
%   \newcolumntype{d}[1]{D{.}{.}{#1}}	% for tables, arg is # decimal places
\usepackage{natbib}			% for author-year citations
   \setlength{\bibhang}{0pt}		% don't use hanging indent format
\usepackage[normalem]{ulem}		% allow underlining
\usepackage{calc}			% calculation package
\usepackage{ifthen}			% control structures
\usepackage{fancyhdr}			% fancy headers and footers
   \fancyhf{}
   \renewcommand{\headrulewidth}{0pt}
   \renewcommand{\footrulewidth}{0pt}
\usepackage{titling}			% more title page control
\usepackage{xspace}			% common-sense spacing after text macro
\usepackage{mdwlist}			% more flexible description lists
\usepackage{color}			% colors and grays
%   \definecolor{subcolor}{rgb}{1.0,0.934,0.801}
   \definecolor{subcolor}{rgb}{0.809,0.852,1.0}
\usepackage{url}			% typesetting urls
   \urlstyle{rm}
%\usepackage[T1]{fontenc}		% for better-looking "_" and tt "{"
\usepackage{ae}				% used instead of above for PDF output
\usepackage{hyperref}			% for hypertext links in PDF
   \hypersetup{letterpaper,plainpages=false,
               pdftitle={The CFD General Notation System Overview and Entry-Level Document},
               pdfauthor={CGNS Project Group},
               colorlinks,
               linkcolor=blue,citecolor=blue,filecolor=red,pagecolor=blue,
               urlcolor=black}
   \renewcommand{\sectionautorefname}{Section}
   \renewcommand{\subsectionautorefname}{Section}
   \renewcommand{\subsubsectionautorefname}{Section}

%
% Page layout
%
\normalsize
\setlength{\oddsidemargin}{0.5in}	% left margins-1.0in for odd/even pages
\setlength{\evensidemargin}{0.0in}
\setlength{\textwidth}{6.0in}		% text width
%\setlength{\marginparwidth}{0in}	% margin note space parameters
%\setlength{\marginparsep}{0in}
%\setlength{\oddsidemargin}{0.25in}	% left margins-1.0in for odd/even pages
%\setlength{\evensidemargin}{0.75in}
%\setlength{\textwidth}{5.5in}		% text width
\setlength{\marginparwidth}{1.1in}	% margin note space parameters
\setlength{\marginparsep}{0.2in}
\reversemarginpar
\setlength{\topmargin}{-0.5in}		% top margin-1.0in
\setlength{\headheight}{0.25in}		% header space parameters
\setlength{\headsep}{0.5in}
\setlength{\textheight}{8.5in}		% text height
\setlength{\topskip}{\baselineskip}	% dist from top of body to 1st baseline
\setlength{\footskip}{0.75in}		% dist from bottom of body to footer
\raggedbottom				% to avoid stretching vertical space
\pagestyle{fancy}			% using fancyhdr package
\fancypagestyle{plain}{%		% page # in corner for plain style
   \fancyhf{}
   \fancyfoot[LE,RO]{\bfseries \thepage}}

%
% Misc spacing parameters
%
%\setlength{\parskip}{\baselineskip}	% blank space between paragraphs
\setlength{\parskip}{0.5\baselineskip}	% blank space between paragraphs
\setlength{\doublerulesep}{0pt}		% for wider \hlines
\setlength{\fboxsep}{5pt}		% for more margin in boxes
\newlength{\saveparindent}		% save for use where LaTeX changes them
\newlength{\saveparskip}		% save for use where LaTeX changes them
\newlength{\savebaselineskip}
\setlength{\saveparindent}{\parindent}
\setlength{\saveparskip}{\parskip}
\setlength{\savebaselineskip}{\baselineskip}
\newlength{\tmplength}			% for temp use wherever needed
\newlength{\tmplengtha}
\newlength{\Pwidth}			% for use in tables
\newlength{\Pwidtha}

%
% User-defined environments
%
% "Definition" list with variable-length entries. Requires
% mdwlist package. (See LaTeX Companion, p 64, and mdwlist documentation.)
\newenvironment{Ventry}[1]%
   {\settowidth{\tmplength}{#1\hspace{1em}}%
    \begin{basedescript}{%
           \desclabelwidth{\tmplength}%
           \desclabelstyle{\nextlinelabel}%
           \renewcommand{\makelabel}[1]{##1\hfil}}}
   {\end{basedescript}}
% "Definition" list with variable-length entries, indented. Requires
% mdwlist package. (See LaTeX Companion, p 64, and mdwlist documentation.)
\newenvironment{Ventryi}[1]%
   {\settowidth{\tmplength}{\hspace{\saveparindent}#1\hspace{1em}}%
    \begin{basedescript}{%
           \desclabelwidth{\tmplength}%
           \desclabelstyle{\nextlinelabel}%
           \renewcommand{\makelabel}[1]{\hspace{\saveparindent}##1\hfil}}}
   {\end{basedescript}}
% Indented definition list with variable-length entries, indented, compact.
%    (See LaTeX Companion, p 64)
\newenvironment{Ventryic}[1]%
   {\begin{list}{}{\renewcommand{\makelabel}[1]{\hspace{\saveparindent}##1\hfil}%
      \settowidth{\labelwidth}{\hspace{\saveparindent}#1\hspace{1em}}%
      \setlength{\leftmargin}{\labelwidth+\labelsep}
      \setlength{\parsep}{0in}%
      \setlength{\itemsep}{0in}}}%
   {\end{list}}
%
% Adapted from comp.text.tex post by Keith Reckdahl
% (reckdahl@am-sparc7.Stanford.EDU)
% Indents text from the left margin by current value of the list
% length \leftmargin
\newenvironment{indleft}%
   {\begin{list}{}
      {\setlength{\topsep}{0pt}%
       \setlength{\listparindent}{0pt}%
       \setlength{\itemindent}{0pt}%
       \setlength{\parsep}{\parskip}%
      }%
      \item[]}%
   {\end{list}}
% As above, but surrounding an alltt environment
\newenvironment{indlefttt}%
   {\begin{list}{}
      {\setlength{\topsep}{0pt}%
       \setlength{\listparindent}{0pt}%
       \setlength{\itemindent}{0pt}%
       \setlength{\parsep}{\parskip}%
      }%
      \item[]%
      \begin{alltt}}%
   {\end{alltt}\end{list}}
% As above, but also with a tabbing environment
\newenvironment{tabindlefttt}%
   {\begin{list}{}
      {\setlength{\topsep}{0pt}%
       \setlength{\listparindent}{0pt}%
       \setlength{\itemindent}{0pt}%
       \setlength{\parsep}{\parskip}%
      }%
      \item[]%
      \begin{alltt}\begin{tabbing}}%
   {\end{tabbing}\end{alltt}\end{list}}

% "changes" environment, to identify changed code by color.  The color is
% given by the argument, with a default of red.  When the environment is
% exited, the color is changed to black.
\newenvironment{changes}[1][red]%
   {\color{#1}}
   {\color{black}}

%
% User-defined commands
%

% Degree symbol
\newcommand{\degree}[1][]{\ensuremath{{#1}^\circ}\xspace}

% Degree Rankine symbol
\newcommand{\degR}{\textsuperscript{$\circ$}\textnormal{R}\xspace}

% Re-def longtable's caption command to use \captionlabelfont from caption pkg
% (basis lifted from longtable.sty)
\makeatletter
\renewcommand{\LT@makecaption}[3]{%
  \LT@mcol\LT@cols c{\hbox to\z@{\hss\parbox[t]\LTcapwidth{%
    \sbox\@tempboxa{{\captionlabelfont #1{#2: }}#3}%
    \ifdim\wd\@tempboxa>\hsize
      {\captionlabelfont #1{#2: }}#3%
    \else
      \hbox to\hsize{\hfil\box\@tempboxa\hfil}%
    \fi
    \endgraf\vskip\baselineskip}%
  \hss}}}
\makeatother

% Make next page odd, with preceding blank page empty (LaTeX Companion, p 93)
\newcommand{\clearemptydoublepage}{\newpage{\pagestyle{empty}\cleardoublepage}}

% Text subscripts, analogous to \textsuperscript command (from comp.text.tex
% post by rf@cl.cam.ac.uk (Robin Fairbairns))
\makeatletter
\DeclareRobustCommand*\tsub[1]{%
  \@tsub{\selectfont#1}}
\def\@tsub#1{%
  {\m@th\ensuremath{_{\mbox{\fontsize\sf@size\z@#1}}}}}
\makeatother

% Text superscripts, but shorter
\newcommand{\tsup}[1]{\textsuperscript{#1}}

% Shortcuts for specific fonts
\newcommand{\bold}[1]{{\normalfont\textbf{#1}}}  % Bold
\newcommand{\ital}[1]{{\normalfont\textit{#1}}}  % Italic
\newcommand{\key} [1]{{\normalfont\texttt{#1}}}  % Fixed font for keywords
\newcommand{\fort}[1]{{\normalfont\texttt{#1}}}  % Fixed font for Fortran stuff

% Shortcuts for math character sizes
\newcommand{\Dstyle}{\displaystyle}
\newcommand{\Tstyle}{\textstyle}
\newcommand{\Sstyle}{\scriptstyle}

% Use section titles as marks (i.e., in headers) instead of subsection titles
\renewcommand{\sectionmark}[1]{\markright{\thesection\ \ #1}}
\renewcommand{\subsectionmark}[1]{}

% "Better" treatment of headings for Appendices  (LaTeX Companion, pp 29,30)
\renewcommand{\appendix}{%
   \newcommand{\app}{%
      \secdef\Appendix\sAppendix}%
   \newcommand{\subapp}{%
      \secdef\subAppendix\subsAppendix}%
   \newcounter{app}%
   \newcounter{subapp}[app]%
   \renewcommand{\theapp}{\Alph{app}}%
   \renewcommand{\thesubapp}{\theapp.\arabic{subapp}}%
 }
\newcommand{\Appendix}[2][?]{%   Complex form
   \refstepcounter{app}%
   \addcontentsline{toc}{section}%
      {\protect\bfseries{\appendixname~\theapp.\ \ } #1}%
   {\noindent\Large\bfseries\appendixname\ \theapp.\ \ #2\par}%
   \sectionmark{#1}\vspace{\baselineskip}}
\newcommand{\sAppendix}[1]{%   Starred form
   \refstepcounter{app}%
   {\noindent\Large\bfseries\appendixname\ \theapp.\ \ #1\par}%
   \vspace{\baselineskip}}

\newcommand{\subAppendix}[2][?]{%   Complex form
   \refstepcounter{subapp}%
   \addcontentsline{toc}{subsection}%
       {\thesubapp\ \ #1}%
   \vspace{\baselineskip}
   {\noindent\large\bfseries\thesubapp\ \ #2\par}%
   \vspace{\baselineskip}}
\newcommand{\subsAppendix}[1]{%   Starred form
   \refstepcounter{subapp}%
   \vspace{\baselineskip}
   {\noindent\large\bfseries\thesubapp\ \ #1\par}%
   \vspace{\baselineskip}}

\begin{document}

\pagenumbering{roman}

%\fancyhead[LE]{\bfseries Overview and Entry-Level Document}
%\fancyhead[RO]{\bfseries \rightmark}
\fancyfoot[LE,RO]{\bfseries \thepage}

%\setlength{\droptitle}{1.0in}
\pretitle{\begin{flushleft}\LARGE%
          \includegraphics[width=2.0in]{../logo/cgns_bw}\\*[0.25in]}
\posttitle{\par\end{flushleft}\vskip 1.0em}
\title{{\bfseries CFD General Notation System\\
Overview and Entry-Level Document}\\*[0.25in]
{\Large Document Version 2.0.20}}
\author{}
\date{}
\maketitle
\thispagestyle{empty}

\clearemptydoublepage
\setlength{\parskip}{0ex}		% remove blank space between paragraphs
\thispagestyle{plain}
\tableofcontents
\setlength{\parskip}{\saveparskip}	% put it back

\clearemptydoublepage
\pagenumbering{arabic}
\setcounter{secnumdepth}{-2}

\fancyhead[LE]{\bfseries Overview and Entry-Level Document}
\fancyhead[RO]{\bfseries \rightmark}

\pdfbookmark[1]{Foreward}{foreward}
\addtocontents{toc}{\protect\contentsline {section}{Foreward}{\thepage}{foreward}}
\hypertarget{foreward}{}
\section*{Foreward}
\thispagestyle{plain}

This document describes the structure of CGNS (CFD General Notation
System), its software, and its documentation.
The Overview, which was written primarily for new and prospective users
of CGNS, (1) introduces terminology, (2) identifies the elements of the
system and their relationships, and (3) describes the various documents
that elaborate the details.
Reading the material on the purpose (\autoref{s:purpose}) and general
description (\autoref{s:description}) of CGNS should help users
determine whether CGNS will meet their needs.
Those wanting a bit more detail should read \autoref{s:elements}, which
describes the various elements of CGNS.
For those interested only in understanding the scope and capabilities of
CGNS, or for the end user unconcerned with the internal workings of the
system, the Overview may prove sufficient documentation by itself.

The Overview also includes certain information that is current as of the
document date but which may change with time.
All information on CGNS compatible ``applications'' software (i.e.,
external programs such as grid generators, flow codes, or postprocessors
in \autoref{s:applications}) is of this type.
Also subject to change is the information on the acquisition of the
software and documentation in \autoref{s:acquiring}, and the current
status of CGNS in \autoref{s:history}.

\setcounter{secnumdepth}{3}

\clearemptydoublepage
\section{Purpose and Scope}
\label{s:purpose}
\thispagestyle{plain}

The general purpose of CGNS is to provide a standard for recording and
recovering computer data associated with the numerical solution of the
equations of fluid dynamics.

CGNS consists of a collection of conventions, and software implementing
those conventions, for the storage and retrieval of CFD (Computational
Fluid Dynamics) data.
The system consists of two parts: (1) a standard format for recording
the data, and (2) software that reads, writes, and modifies data in that
format.
The format is a conceptual entity established by the documentation, and
is intended to be general, portable, expandable, and durable.
The software is a physical product supplied to enable developers to
access and produce data recorded in that format.
All CGNS software is completely free and open to anyone. 

The CGNS standard, applied through the use of the supplied software, is
intended to:

\begin{itemize*}
   \item facilitate the exchange of CFD data
   \begin{itemize*}
      \item between sites
      \item between applications codes
      \item across computing platforms
   \end{itemize*}
   \item stabilize the archiving of CFD data
\end{itemize*}

The principal target of CGNS is data normally associated with
compressible viscous flow (i.e., the Navier-Stokes equations), but the
standard is also applicable to subclasses such as Euler and potential
flows.
The CGNS standard addresses the following types of data.

\begin{itemize*}
\item Structured, unstructured, and hybrid grids
\item Flow solution data, which may be nodal, cell-centered,
      face-centered, or edge-centered
\item Multizone interface connectivity, both abutting and overset
\item Boundary conditions
\item Flow equation descriptions, including the equation of state,
      viscosity and thermal conductivity models, turbulence models,
      and multi-species chemistry models
\item Time-dependent flow, including moving and deforming grids
\item Dimensional units and nondimensionalization information
\item Reference states
\item Convergence history
\item Association to CAD geometry definitions
\end{itemize*}

Much of the standard and the software is applicable to computational
field physics in general.
Disciplines other than fluid dynamics would need to augment the data
definitions and storage conventions, but the fundamental database
software, which provides platform independence, is not specific to fluid
dynamics.


\clearemptydoublepage
\section{General Description}
\label{s:description}
\thispagestyle{plain}

A CGNS database describes the current state of one or more entire CFD
(Computational Fluid Dynamics) problems, including the following:

\begin{itemize*}
   \item grid
   \item flowfield
   \item boundary conditions
   \item topological connection information
   \item auxiliary data (e.g., nondimensionalization parameters,
         reference states)
\end{itemize*}

Not all of these data need to be present at any particular time.
The overall view is that of a shared database that can be accessed
by the various software tools common to CFD, such as solvers, grid
generators, field visualizers, and postprocessors.
Each of these ``applications'' serves as an editor of the data, adding to,
modifying, or interpreting it according to that application's specific
role.

CGNS conventions and software provide for the recording of a complete
and flexible problem description.
The exact meaning of a subsonic inflow boundary condition, for example,
can be described in complete detail if desired.
User comments can be included nearly anywhere, affording the
opportunity, for instance, for date stamping or history information to
be included.
Dimension and sizing information is carefully defined.
Any number of flow variables may be recorded, with or without standard
names, and it is also possible to add user-defined or site-specific
data.
These features afford the opportunity for applications to perform
extensive error checking if desired.

Because of this generality, CGNS provides for the recording of much more
descriptive information than current applications normally use.
However, the provisions for this data are layered so that much of it is
optional.
It should be practical to convert most current applications to CGNS with
little or no conceptual change, retaining the option to take advantage
of more detailed descriptions as that becomes desirable.

CGNS specifications currently cover the bulk of CFD data that one might
wish to exchange among sites or applications; for instance, nearly any
type of field data can be recorded, and, based on its name, found and
understood by any code that needs it.
Global data (e.g., freestream Mach number, Reynolds number, angle
of attack) and physical modeling instructions (e.g., thin layer
assumptions, turbulence model) may be specified.
Nevertheless, there are items specific to individual applications for
which there is currently no specification within CGNS.
Most commonly, these are operational instructions, such as number of
sweeps, solution method, multigrid directives, and so on.
Owing to the miscellaneous nature of this data, there has been no
attempt to codify it within a global standard.
It is therefore expected that many applications will continue to require
small user-generated input files, presumably in ASCII format.

CGNS itself does not initiate action or undertake any function normally
handled by the operating system.
The user still performs CFD tasks according to existing processes.
This includes selecting the computing platform, maintaining the files,
and launching the applications.

However, the ease of communication between applications that CGNS
provides should motivate the development of new batch and interactive
mechanisms for the convenient application of CFD tools.


\clearemptydoublepage
\section{Elements and Documentation}
\label{s:elements}
\thispagestyle{plain}

\subsection{Introduction}

CGNS concerns itself with the recording and retrieval of data associated
with the computation of fluid flows.
Included are such structures as grids, flowfields, boundary conditions,
and zone connectivity information.
CGNS ``understands'' this data in the sense that it contains conventions
for recording it based on its structure and its role in CFD.

The underlying design of CGNS is that of a common database that is
accessed (read, written, or modified) by a collection of ``applications''
programs such as solvers, grid generators, and postprocessors.

CGNS itself does not synthesize, modify, or interpret the data it
stores.
The applications create, edit, or display the data; CGNS is limited to
recording and retrieving it.
Each application's program accesses the data directly using CGNS
function calls installed in the application by its developer.
The applications are not regarded as part of CGNS itself.

CGNS is passive.
It does not initiate action and cannot ``push'' information into the
applications codes or ``pull'' information out.
Rather, the codes must request the information they seek and store the
information they produce.
The applications must be launched by a user who organizes the location
and content of the database.
The process and sequence of events remain under user control.
Thus CGNS facilitates, but does not incorporate, the development of
batch or interactive environments designed to control the CFD process.

The elements of CGNS address all activities associated with the
storage of the data on external media and its movement to and from the
applications programs.
These elements include the following:

\begin{itemize}
\item The ADF (``Advanced Data Format'') Database Manager, consisting
      of both a file format specification and its I/O software, which
      handles the actual reading and writing of data from and to
      external storage media.
      In versions 2.4 and later of the CGNS software, the user can
      choose to use the HDF5 storage file format (supported externally
      by NCSA) in place of ADF.
\item The Standard Interface Data Structures (SIDS), which specify the
      intellectual content of CFD data and the conventions that govern
      naming and terminology.
\item The SIDS-to-ADF File Mapping and SIDS-to-HDF File Mapping
      conventions, which specify the exact location where the CFD
      data defined by the SIDS is to be stored within the ADF or HDF5
      file(s).
\end{itemize}

It should be noted that because of HDF5's parallel and compression      
capability as well as its support, the CGNS Steering Committee has made 
the decision to slowly transition (beginning in 2005) to HDF5 as the    
official data storage mechanism.                                        
However, ADF will continue to be available for use, with the CGNS
mid-level library capable of (1) using either format and (2) translating
back and forth between the two.

The ADF and HDF5 Database Managers provide minimal collections of data
access functions.
This collection we refer to as the ``ADF Core,'' and the ``HDF5 Core.''
In principle, it is possible to install CGNS I/O into an application
using only the ADF Core or HDF5 Core routines.
However, such an approach would require the installer to access the data
at a very fundamental level and would result in lengthy sequences of
core function calls.
Therefore, the system also includes a Mid-Level Library, an API
(Application Programming Interface) that contains additional routines
intended to facilitate higher-level access to the data.
These are CFD-knowledgeable routines suitable for direct installation
into applications codes.

The following sections discuss in more detail the roles of the CGNS
elements and introduce their documentation.

\subsection{ADF Files and the ADF Core}

In this section, we refer solely to the ADF Database Manager.  
However, the HDF5 storage format, which can be invoked as an alternative
storage method in place of ADF, possesses similar functionality.
Refer to the HDF5 website at \url{http://hdf.ncsa.uiuc.edu/HDF5/} for
more information.

ADF (Advanced Data Format), with its software implementation called the
ADF Core, constitutes an extremely general database manager that is
particularly suited to the storage of large numerical data sets.
The Core is composed of a set of routines that perform standard
operations on a database consisting of ADF files.
Typical operations include the following: open, close, create, delete,
read, and write.

An ADF file is, by definition, a file that has been written by the ADF
Core.
ADF files are written onto the external medium in a binary format that
is completely determined by the Core.
It is by this means that file portability is ensured.
Consequently, the file contents are not accessible except by using the
Core routines.
The Core routines are internally documented, but they contain no
user-serviceable parts and \emph{should not be modified}.

The principal advantages of such an independently designed storage
system are portability, data integrity, and the opportunity for
self-identification of files.
The ADF Core can always open and read an ADF file regardless of its
origin or specific content.
The universal availability of ADF, resulting from its nonproprietary
nature, is also an important feature.

All ADF Core routines are coded in C, but both C and Fortran function
calls are provided.
The Core routines provide only the minimum functionality required to
access the database.

The conceptual form of an ADF file is that of a tree, each node of which
has the same internal structure.
Each node contains the following:

\begin{itemize*}
\item identification information
\item pointers to its subnodes
\item descriptors of any data recorded at the node
\item the data itself
\end{itemize*}

This format was chosen with the storage of scientific data in mind, but
it is general enough that the ADF system could be used to store a wide
variety of types of information.
There is nothing in the ADF file specifications that restricts the data
types to CFD-related information.

It is possible, through a mechanism of links, for a database implemented
in ADF to span more than one ADF file.
In the case of CFD, for instance, this feature allows multiple flow
solutions to share the same grid without duplicating the data.

ADF contains no mechanisms to ensure the security or availability of ADF
files.
It is the user's responsibility to keep track of the locations, file
permissions, and names of ADF files.
ADF does not assume or modify any functions of the operating system.
In particular, there is no file ``locking''; i.e., there is nothing to
prevent two applications from accessing the same file simultaneously.
ADF file links are ``soft'' --- if the linked file is moved or deleted, the
connection is broken.

A more thorough overview of ADF may be found in the \textit{SIDS-to-ADF
File Mapping} document, and complete documentation appears in the
\textit{ADF User's Guide}.

\subsection{Standard Interface Data Structures (SIDS)}

The establishment of a standard for storing CFD-related information
requires a detailed specification of the content and meaning of the data
to be stored.
For example, it is necessary to state the meaning of the words ``boundary
condition'' in a form sufficiently concrete to be recorded precisely, and
yet sufficiently flexible to embrace current and future practice.
The \textit{Standard Interface Data Structures}
(SIDS) document describes this ``intellectual content'' of CFD-related
data in detail.

An exact description of the intellectual content is required not only
to define the precise form of the data but also to guarantee that the
meaning of the data is consistently interpreted by practitioners.
Thus the SIDS include a collection of naming conventions that
specify the precise meaning of nomenclature (e.g., the strings
\fort{DensityStatic} and \fort{BCWallViscous}).

The SIDS are written in a self-contained C-like descriptive language.
SIDS data structures are defined in a hierarchical manner in which more
complex entities are built from simpler ones.
These structures are closely reflected in CGNS-compliant ADF (or HDF5)
files: simple entities are often stored in single nodes, while more
complex structures are stored in entire subtrees.

\subsection{SIDS-to-ADF File Mapping and SIDS-to-HDF File Mapping}

\subsubsection{General File Mapping Principles}

Because of the generality of the ADF (or HDF5) tree structure, there are
many conceivable means of encoding CFD data.
But for any application to access, say, the boundary conditions for zone
``UnderWing'', requires a single convention with regard to where in the
file that data has been stored.
The \textit{SIDS-to-ADF File Mapping} document, sometimes referred to
as the ``File Mapping,'' establishes the precise ADF node, and location
within that node, where each piece of CGNS data should be recorded.
Similarly, the SIDS-to-HDF File Mapping document does the same thing for
HDF5.
The CGNS Mid-Level Library relies on the File Mapping to locate
CFD-related data within the file.

The mapping provides ADF (or HDF5) locations for an extensive set of CFD
data.
Most applications will make use of only a small subset of this data.
Further, inasmuch as applications are viewed as editors that are in the
process of building the database, most of them are intended for use on
incomplete data sets.
Therefore, it is not required that all the data elements specified by
the CGNS conventions be complete in order for a database to be CGNS
compliant.
The user must ensure that the current state of the database will support
whatever application he may launch.
Of course, the application should gracefully handle any absence or
deficiency of data.

CGNS conventions do not specify the following:

\begin{itemize*}
\item the use the applications programs may make of the data
\item the means by which the applications programs modify the data
\item the form in which the data is stored internal to an application
\end{itemize*}

\noindent
The validity, accuracy and completeness of the data are determined
entirely by the applications software.

The tree structure of ADF (or HDF5) makes it possible for applications
to ignore data for which they have no use.
(In fact, they cannot even discover the data's existence without a
specific inquiry.)
Therefore, it is permissible for an ADF (or HDF5) file containing a CGNS
database to contain additional nodes not specified by the mapping.
Such nodes will be disregarded by software not prepared to use them.
However, if data essential to the CFD process is stored in a manner
not consistent with CGNS conventions, that data becomes invisible and
therefore useless to other applications.

Note that the SIDS serve not only to facilitate the mapping of data onto
the ADF (or HDF5) file structure but also to standardize the meaning of
the recorded data.
Thus there are two kinds of conventions operative within CGNS.
Adherence to the File Mapping conventions guarantees that the software
will be able to find and read the data.
Adherence to the SIDS guarantees uniformity of meaning among users and
between applications.
The \textit{SIDS-to-ADF File Mapping} document establishes the
use of ADF in the context of CGNS (the \textit{SIDS-to-HDF File
Mapping} document does the same thing for HDF5); the SIDS define the
nomenclature, content, and meaning of the stored data.

The File Mapping generally avoids the storage of redundant data.
Sometimes an application may require an alternate (but intellectually
equivalent) form of the data; in such cases it is recommended that the
alternate form be prepared at the time of use and kept separate from the
CGNS data.
This avoids habitual reliance on the alternate form, which would
invalidate the standard.
If the alternate form is appended to the ADF (or HDF5) file, care must
be taken to update the primary (CGNS) form whenever permanent changes
are made.

\subsubsection{Structure of ADF Files and ADF Nodes}

In this section, the structure of ADF files and nodes is discussed.
The HDF5 storage format, which can be invoked as an alternative storage
method in place of ADF, is not referred to at all in this section.
Refer to the HDF5 website at \url{http://hdf.ncsa.uiuc.edu/HDF5/} for
more information on the structure of HDF5 files and nodes.

Further description of the File Mapping requires some understanding of
the internal structure of ADF Files.
In this document, we summarize this structure as it is used in CGNS.
The \textit{SIDS-to-ADF File Mapping} document contains a complete
description of the ADF file structure as used in CGNS.
The \textit{ADF User's Guide} contains a complete description of ADF
file structure without regard to its use in CGNS.

An ADF file consists of a collection of elements called nodes.
These nodes are arranged in a tree structure that is logically similar
to a UNIX file system.
The nodes are said to be connected in a ``child-parent'' relationship
according to the following simple rules:

\begin{enumerate}
\item Each node may have any number of child nodes.
\item Except for one node, called the root, each node is the child of
      exactly one other node, called its parent.
\item The root node has no parent.
\end{enumerate}

Each node in an ADF file has exactly the same internal structure.
The entities associated with each node are the following:

\begin{itemize*}
\item Node Identifier (ID)
\item Name
\item Label
\item Data Type
\item Dimension
\item Dimension Values
\item Data
\item Child Table
\end{itemize*}

\paragraph{Node Identifier.}
The Node ID is a floating point number assigned by the system when the
ADF file is opened or created.
Applications may record the ID and use it to return directly to the
corresponding node when required.
The Node ID is valid only while the file is open; subsequent openings of
the same file may be expected to yield different IDs.

\paragraph{Name.}
The Name field holds a character string chosen by the user or specified
by the SIDS to identify the particular instance of the data being
recorded.

\paragraph{Label.}
The Label, also a character string, is specified by the CGNS mapping
conventions and identifies the kind of data being recorded.
For example, a node with label \fort{Zone\_t} may record (at and below
it) information on the zone with Name ``UnderWing.''
It is a general (ADF) rule that no node may have more than one child
with the same name, but the CGNS mapping conventions commonly specify
many children with the same label.
For some nodes, the mapping conventions specify that the name
field has significance for the meaning of the data (e.g.,
\fort{EnthalpyStagnation}).
Although the user may specify another name, these ``paper'' conventions
serve the transfer of data between users and between applications.
These names and their meanings are established by the SIDS.

\paragraph{Data Type, Dimension, Dimension Values, Data.}
ADF nodes may or may not contain data.
For those that do, CGNS specifies a single array whose type (integer,
etc.), dimension, and size are recorded in the Data Type, Dimension, and
Dimension Value fields, respectively.
The mapping conventions specify some nodes that serve to establish the
tree structure and point to further data below but contain no data
themselves.
For these nodes, the Data Type is \fort{MT}, and the other fields are
empty.

\paragraph{Child Table.}
The Child Table contains a list of the node's children.
It is maintained by the system as children are created and deleted.
Its role in CGNS is the same as in ADF in general.

\subsubsection{High-Level Organization of the CGNS Database}

For a full specification of the location of CFD data in the ADF (or
HDF5) file, the user should see the \textit{SIDS-to-ADF File Mapping}
document or the \textit{SIDS-to-HDF File Mapping} document.
For convenience, we summarize the high-level structure below.

A CGNS database consists of a tree of nodes implemented as all or part
of one or more ADF (or HDF5) files.
All information is identified by and accessed through a single node in
one of these files.
There is thus no notion of a CGNS \textit{file}, only of a CGNS
\textit{database} implemented within one or more \textit{ADF (or HDF5)
files}.

By definition, the root node of a CGNS database has the Label
\fort{CGNSBase\_t}.
The Name of the database can be specified by the user and is stored in
the ``Name'' field of the \fort{CGNSBase\_t} node.
Note that this is the root of the CGNS database, not the root node of
the ADF (or HDF5) file.
Current CGNS conventions require that the \fort{CGNSBase\_t} node be
located directly below the ADF (or HDF5) root node.

An ADF (or HDF5) file may contain multiple CGNS databases, and thus
multiple \texttt{CGNSBase\_t} nodes.
However, each node labeled \fort{CGNSBase\_t} in a single ADF (or HDF5)
file must have a unique name.
The user or application must know the name of the file containing
the entry-level node and, if there is more than one node labeled
\fort{CGNSBase\_t} in that file, the name of the database as well.

Below the \fort{CGNSBase\_t} node, the mapping conventions specify a
subnode for each zone.
This node has label \fort{Zone\_t}.
Its Name refers to the particular zone whose characteristics are
recorded at and below the node, such as ``UnderWing.''
In general, names can be specified by the user, but defaults are
specified for nodes that the user does not choose to name.
For the \fort{Zone\_t} nodes, the defaults are \fort{Zone1},
\fort{Zone2}, and so forth, in order of creation.
A similar convention for default names applies elsewhere.
It is impossible to create a node without a name (or with a name of zero
length).
The CGNS Mid-Level Library conforms to the default convention.

Below each zone node will be found nodes for the grid, flowfield,
boundary conditions, and connectivity information; these, in turn, are
parents of nodes specifying extent, spatial location, and so on.

The file mapping specifies that one or more ``Descriptor'' nodes may be
inserted anywhere in the file.
Descriptor nodes are used to record textual information regarding the
file contents.
The size of Descriptor nodes is unlimited, so entire documents could be
named and stored within the data field if desired.
Descriptors are intended to store human-readable text, and they are not
processed by any supplied CGNS software (except, of course, that the
text may be stored and retrieved).

It is possible, by using the linking capability of CGNS, for a child of
any node to be a node in another ADF (or HDF5) file, or elsewhere within
the same file.
This mechanism enables one database to share a grid, for example, with
another database without duplicating the information.

\subsection{Mid-Level Library, or API}

The CGNS Mid-Level Library, or Applications Programming Interface
(API), is one of the most directly visible parts of CGNS, and it is of
particular interest to applications code developers.
It consists of a set of routines that are designed to allow applications
to access CGNS data according to the role of the data in CFD.
Unlike the ADF (or HDF5) Core, routines in the CGNS Mid-Level Library
``understand'' the SIDS-defined CFD data structures and the File
Mapping.
This enables applications developers to insert CGNS I/O into their
applications without having detailed knowledge of the File Mapping.
For instance, an application might use CGNS mid-level calls to retrieve
all boundary conditions for a given zone.

The \textit{CGNS Mid-Level Library} document contains complete
descriptions and usage instructions for all mid-level routines.
As with the ADF (or HDF5) Core, all calls are provided in both C and
Fortran.

\subsection{Documentation}

The five CGNS elements described above are documented individually,
as follows:

\begin{itemize*}
\item \textit{ADF User's Guide}
\item \textit{Standard Interface Data Structures}
\item \textit{SIDS-to-ADF File Mapping Manual}
\item \textit{SIDS-to-HDF File Mapping Manual}
\item \textit{Mid-Level Library}
\end{itemize*}

\noindent
In addition, the following documentation is also recommended:

\begin{itemize*}
\item \textit{CGNS Overview and Entry-Level Document} (this document)
\item ``The CGNS System'', AIAA Paper 98-3007
\item ``Advances in the CGNS Database Standard for Aerodynamics and
      CFD'', AIAA Paper 2000-0681
\item ``CFD General Notation System (CGNS): Status and Future
      Directions'', AIAA Paper 2002-0752
\end{itemize*}

\noindent
The specific documents of interest vary with the level of intended use
of CGNS.

\subsubsection{Prospective Users}

Prospective users are presumably unfamiliar with CGNS.
They will probably wish to begin with the current Overview document, or,
if they require more detailed information, the AIAA papers listed above.
Beyond that, most will find a quick read of the \textit{SIDS-to-ADF
File Mapping Manual} (or the \textit{SIDS-to-HDF File Mapping Manual})
enlightening as to the logical form of the contents of CGNS files.
Browsing the figures in the File
Mapping Manual, as well as the SIDS
itself, will provide some feel for the scope of the system.
The \textit{User's Guide to CGNS}, and
the \textit{CGNS Mid-Level Library}
document, should give an indication of what might be required to
implement CGNS in a given application.
Prospective users should probably not concern themselves with the
details of ADF (or HDF5).

\subsubsection{End Users}

The end user is the practitioner of CFD who generates the grids, runs
the flow codes and/or analyzes the results.
For this user, a scan of this Overview document will sufficiently
explain the overall workings of the system.
This includes end user responsibilities for matters not governed by
CGNS, such as the maintenance of files and directories.
The end user will also find useful the
\textit{User's Guide to CGNS}, as
well as those portions of the SIDS
which deal with standard data names.
The AIAA papers listed above may also be useful if more details about
the capabilities of CGNS are desired.

\subsubsection{Applications Code Developers}

The applications code developer builds or maintains code to support the
various sub-processes encountered in CFD, e.g., grid generation, flow
solution, post-processing, or flow visualization.
The code developer must be able to install CGNS compliant I/O.
The most convenient method for doing so is to utilize the CGNS Mid-Level
Library.
The \textit{User's Guide to CGNS} is the
starting point for learning to use the Mid-Level Library to create and
use CGNS files.
The \textit{CGNS Mid-Level Library}
document itself should also be considered essential.
This library of routines will perform the most common I/O operations in
a CGNS-compliant manner.
However, even when the Mid-Level Library suffices to implement all
necessary I/O, an understanding of the
file mapping and SIDS will be useful.
It will likely be necessary to consult the
SIDS to determine the precise meaning
of the nomenclature.

Applications code developers wishing to read or write data in ADF (or
HDF5) files, that isn't supported by the Mid-Level Library, will have
to use the ADF (or HDF5) Core routines directly, and will need the
\textit{ADF User's Guide} for ADF, and separate HDF5 documentation (not
available here) for HDF5.

\subsubsection{CGNS System Developers}

CGNS System development can be kept somewhat compartmentalized.
Developers responsible for maintenance of the ADF Core, or the building
of supplements to the ADF Core, need not concern themselves with
documentation other than the \textit{ADF User Guide}.
(Development and maintenance of HDF5 is under the purview of NCSA, so
has no relevance here.)
System developers wishing to add to the CGNS Mid-Level Library will need
all the documents.
Theoretical developments, such as extensions to the SIDS, may possibly
be undertaken with a knowledge of the SIDS alone, but such contributions
must also be added to the File Mapping before they can be implemented.


\clearemptydoublepage
\section{Applications Software}
\label{s:applications}
\thispagestyle{plain}

The development of CGNS-compliant applications, e.g., grid generators,
postprocessors, and the like, has not been a direct undertaking of the
CGNS team.
Rather, it has been the intent to make the attractiveness of
interoperable CFD applications, together with general acceptance of
the CGNS standard by Boeing, NASA, and others, sufficient to induce
applications developers to incorporate CGNS I/O into their offerings.

Several CGNS-compatible applications have indeed been developed, and
more continue to appear.
This section lists the known applications at the time of publication.
Please contact Tony Iannetti (\textit{Anthony.C.Iannetti@nasa.gov})
to modify or add to the lists.

\subsection{CGNS-Compatible Applications}

This section lists major software packages that read and/or write
CGNS files.
The software listed here is taken from Table 1 of AIAA Paper 2002-0752,
and is for informational purposes only.
The CGNS Project Group cannot endorse specific software packages.

\setlength{\LTleft}{\fill}
\setlength{\LTright}{\fill}
\setlength{\LTcapwidth}{5.0in}
\settowidth{\tmplength}{Advanced Visual Sys.}
\settowidth{\tmplengtha}{ANSYS ICEM CFD}
\setlength{\Pwidth}{\linewidth-4\tabcolsep-\tmplength-\tmplengtha}
\begin{longtable}{@{}l l >{\itshape}p{\Pwidth}@{}}
\caption[CGNS-Compatible Grid Generators and Pre-Processors]{\textbf{CGNS-Compatible Grid Generators and Pre-Processors}}
\label{t:gridgen}
\\ \hline\hline \\*[-2ex]
\bold{Organization} & \bold{Application} & \bold{Contact}
\\*[1ex] \hline\hline \\*[-2ex]
\endfirsthead

\multicolumn{3}{l}{{\bfseries\autoref{t:gridgen}: CGNS-Compatible Grid Generators} (\emph{Continued})}
\\*[1ex] \hline\hline \\*[-2ex]
\bold{Organization} & \bold{Application} & \bold{Contact}
\\*[1ex] \hline\hline \\*[-2ex]
\endhead

\\*[-2ex]\hline
\multicolumn{3}{r}{\emph{Continued on next page}} \\
\endfoot
\\*[-2ex] \hline\hline
\endlastfoot
\textit{Commercial} \\
Advanced Visual Sys. & ANSYS ICEM CFD & \url{http://cfl3d.larc.nasa.gov/Cfl3dv6/cfl3dv6.html#contact} \kill
Ansys, Inc.     & ANSYS ICEM CFD & \url{ansysinfo@ansys.com} \\
NUMECA          & IGG           & \url{http://www.numeca.be/index.php?id=offices} \\
Pointwise, Inc. & GRIDGEN       & \url{http://www.pointwise.com/pointwise/contact.shtml} \\
\\
\textit{Industry} \\
Boeing Rocketdyne        & APPT   & \\
\end{longtable}

\enlargethispage{2\baselineskip}

\setlength{\LTleft}{\fill}
\setlength{\LTright}{\fill}
\begin{longtable}{@{}l l >{\itshape}p{\Pwidth}@{}}
\caption[CGNS-Compatible Flow Solvers]{\textbf{CGNS-Compatible Flow Solvers}}
\label{t:solvers}
\\ \hline\hline \\*[-2ex]
\bold{Organization} & \bold{Application} & \bold{Contact}
\\*[1ex] \hline\hline \\*[-2ex]
\endfirsthead

\multicolumn{3}{l}{{\bfseries\autoref{t:solvers}: CGNS-Compatible Flow Solvers} (\emph{Continued})}
\\*[1ex] \hline\hline \\*[-2ex]
\bold{Organization} & \bold{Application} & \bold{Contact}
\\*[1ex] \hline\hline \\*[-2ex]
\endhead

\\*[-2ex]\hline
\multicolumn{3}{r}{\emph{Continued on next page}} \\
\endfoot
\\*[-2ex] \hline\hline
\endlastfoot
\textit{Commercial} \\
Advanced Visual Sys. & ANSYS ICEM CFD & \url{http://cfl3d.larc.nasa.gov/Cfl3dv6/cfl3dv6.html#contact} \kill
Ansys, Inc.    & CFX5         & \url{ansysinfo@ansys.com} \\
               & CFX-TASCFlow & \url{ansysinfo@ansys.com} \\
Aerosoft, Inc. & GASP         & \url{http://www.aerosft.com/Misc/contact.php} \\
CD ADAPCO      & STAR-CD      & \url{http://www.cd-adapco.com/about/global_locations.html} \\
Fluent, Inc.   & Fluent       & \url{http://www.fluent.com/contact/} \\
Newmerical Tech. & FENSAP     & \url{http://www.newmerical.com/public/eng/index.php?pageId=about&language=en} \\
NUMECA         & EURANUS      & \url{http://www.numeca.be/index.php?id=offices} \\
\newpage
\textit{Government} \\
DLR/MTU        & TRACE    & \url{Frank.Eulitz@dlr.de} \\
NASA Langley   & PEGSUS   & \url{Pieter.G.Buning@nasa.gov} \\
               & OVERFLOW & \url{Pieter.G.Buning@nasa.gov} \\
               & Cart3D   & \url{http://people.nas.nasa.gov/~aftosmis/} \\
NASA Glenn     & NCC      & \url{Nan-Suey.Liu-1@nasa.gov} \\
NASA Langley   & CFL3D    & \url{http://cfl3d.larc.nasa.gov/Cfl3dv6/cfl3dv6.html#contact} \\
               & UPS      & \\
NPARC Alliance & WIND     & \url{npsupport@info.arnold.af.mil} \\
ONERA          & elsA     & \url{http://elsa.onera.fr/elsA/contacts.html} \\
\\
\textit{Industry} \\
Boeing Seattle           & TLNS3D & \\
Rolls-Royce Allison      & ADPAC  & \\
Rolls-Royce Oxford       & HYDRA  & \url{pierre.moinier@comlab.ox.ac.uk}
\end{longtable}

\setlength{\LTleft}{\fill}
\setlength{\LTright}{\fill}
\begin{longtable}{@{}l l >{\itshape}p{\Pwidth}@{}}
\caption[CGNS-Compatible Post Processors]{\textbf{CGNS-Compatible Post Processors}}
\label{t:postprocess}
\\ \hline\hline \\*[-2ex]
\bold{Organization} & \bold{Application} & \bold{Contact}
\\*[1ex] \hline\hline \\*[-2ex]
\endfirsthead

\multicolumn{3}{l}{{\bfseries\autoref{t:postprocess}: CGNS-Compatible Post Processors} (\emph{Continued})}
\\*[1ex] \hline\hline \\*[-2ex]
\bold{Organization} & \bold{Application} & \bold{Contact}
\\*[1ex] \hline\hline \\*[-2ex]
\endhead

\\*[-2ex]\hline
\multicolumn{3}{r}{\emph{Continued on next page}} \\
\endfoot
\\*[-2ex] \hline\hline
\endlastfoot
\textit{Commercial} \\
Advanced Visual Sys. & ANSYS ICEM CFD & \url{http://cfl3d.larc.nasa.gov/Cfl3dv6/cfl3dv6.html#contact} \kill
Advanced Visual Sys. & AVS/Express  & \url{http://www.avs.com/contact/} \\
Aerosoft, Inc.       & GASP         & \url{http://www.aerosft.com/Misc/contact.php} \\
Tecplot, Inc.        & Tecplot      & \url{http://www.tecplot.com/contact_us/} \\
CEI                  & Ensight      & \url{http://www.ensight.com/contact-us-3.html} \\
                     & Ensight Gold & \url{http://www.ensight.com/contact-us-3.html} \\
Fluent, Inc.         & Fluent       & \url{http://www.fluent.com/contact/} \\
Ansys, Inc.          & ANSYS ICEM CFD & \url{ansysinfo@ansys.com} \\
Intelligent Light    & Fieldview    & \url{fieldview@ilight.com} \\
Newmerical Tech.     & DROP3D       & \url{http://www.newmerical.com/public/eng/index.php?pageId=about&language=en} \\
                     & ICE3D        & \url{http://www.newmerical.com/public/eng/index.php?pageId=about&language=en} \\
NUMECA               & CFView       & \url{http://www.numeca.be/index.php?id=offices} \\
\\
\textit{Government} \\
NASA Ames            & PLOT3D       & \url{Cetin.C.Kiris@nasa.gov}
\end{longtable}

\newpage
\subsection{CGNS Utilities}

A variety of CGNS utility packages have been developed, both by CGNS
team members and by CGNS users.
While not formally a part of CGNS itself, these utilities are available
from SourceForge, at \textit{http://sourceforge.net/projects/cgns},
and/or from the CGNS web site, at \textit{http://cgns.sourceforge.net/}.

\begin{Ventryi}{Plot3dg\_to\_CGNS}
\item [Plot3dg\_to\_CGNS]
      Converts a ``standard'' Plot3D grid file to a CGNS grid file.\\
      Contact: Chris Rumsey (\textit{Christopher.L.Rumsey@nasa.gov})
\item [CGNS\_to\_Plot3d]
      Converts a CGNS file with a cell-centered solution to ``standard''
      Plot3D grid and solution files.\\
      Contact: Chris Rumsey (\textit{Christopher.L.Rumsey@nasa.gov})
\item [CGNS\_readhist]
      Reads a CGNS file and writes history data to a formatted file.\\
      Contact: Chris Rumsey (\textit{Christopher.L.Rumsey@nasa.gov})
\item [FTU]
      Converts CGNS files to and from Plot3D format.
      Has a text-based menu allowing the manipulation of a CGNS base.\\
      Contact: Bob Bush (\textit{bushrh@pweh.com})
\item [CGNSview]
      An CGNS file viewer and editor, written in \textit{Tcl/Tk}.
      The GUI allows complete editing of CGNS files (both ADF and HDF5),
      and handles links and data conversions. It works at the CGIO
      interface level, and thus is not affected by ``bad'' CGNS files.
      It includes built-in CGNS node and label information, and
      documentation support through HTML (including access to all
      CGNS docs).\\
      Contact: Bruce Wedan (\textit{brucewedan@gmail.com})
\item [CGNS Tools]
      A variety of tools for viewing, editing, and manipulating the
      contents of CGNS files; requires \textit{Tcl/Tk} 8.3 or above.
      Included are:
      \begin{itemize*}
      \item \textit{CGNSplot}, for plotting grids
      \item \textit{CGNScalc}, a calculator using data in CGNS files
      \item Utilities for converting from Patran, and from/to PLOT3D and
            Tecplot
      \item Utilities for modifying the grid location
            (vertex or cell-center), solution variables
            (primitive or conservative), and data class
            (\fort{Dimensional}, \fort{NormalizedByDimensional}, or
            \fort{NormalizedByUnknownDimensional})
      \item Utilities for creating a subset or interpolated version of
            a CGNS file
      \end{itemize*}
      These may be used as stand-alone utilities, or initiated from
      \textit{CGNSview}, which is included.\\
      Contact: Bruce Wedan (\textit{brucewedan@gmail.com})
\item [CGNS Viewer]
      An alpha release of another ADF/CGNS file viewer and editor.
      Requires the GTK+ library version 1.2 or newer (installed by default
      on newer Linux systems using Gnome; otherwise found at
      {\itshape\url{http://www.gtk.org/}}).\\
      Contact: Christian Lundh (\textit{xistan@yahoo.se})
\item [PyCGNS]
      A Python binding to the ADF Core and Mid-Level Library.
      For more information, see the pyCGNS home page at
      {\itshape\url{http://elsa.onera.fr/CGNS/releases/}}.\\
      Contact: Marc Poinot (\textit{Marc.Poinot@onera.fr})
\item [ADFM]
      A development version of software for managing an in-memory
      representation of ADF trees.
      For more information, see ``Proposal and Prototype for an in-memory
      representation of ADF trees'' at
      {\itshape\url{http://cgns.sourceforge.net/utilities/ADFM/ADFM.pdf}}.\\
      Contact: Marc Poinot (\textit{Marc.Poinot@onera.fr})
\item [CGNS++]
      A C++ binding to the ADF Core and Mid-Level Library.
      For more information, see the CGNS++ home page at
      \textit{http://cgnspp.sourceforge.net/}.\\
      Contact: Manuel Kessler (\textit{kessler@iag.uni-stuttgart.de})
\end{Ventryi}


\clearemptydoublepage
\section{Acquiring CGNS}
\label{s:acquiring}
\thispagestyle{plain}

\subsection{Software}

The CGNS software is available free of charge, under the terms of the
CGNS License, from SourceForge, at
\textit{http://sourceforge.net/projects/cgns}.
Also available there are the \textit{cgnstools} utilities, the source
code examples from \textit{A User's Guide to CGNS}, and additional
Fortran source code examples.

The CGNS Library contains source code for the Mid-Level Library,
the CGIO core, and the ADF and HDF5 implementations,
plus CMake and configure scripts for building the
library for a variety of platforms.

\subsection{Documentation}

The CGNS documentation may be accessed via the CGNS Documentation home
page at {\itshape\url{http://www.grc.nasa.gov/WWW/cgns/}}.
Documentation may also be available for the current beta version of CGNS.
All the CGNS documentation is available in both PDF and HTML forms.

In addition to the CGNS documentation itself, several conference papers
and slide presentations are available, as well as minutes from the CGNS
meetings and telecons.


\clearemptydoublepage
\section{History and Current Status}
\label{s:history}
\thispagestyle{plain}

\subsection{Initial Development}

\enlargethispage{\baselineskip}
The initial development of CGNS took place from March 1995 through
May 1998, when Version 1.0 of the CGNS software was released.
The project originated through a series of meetings in 1994-1995 between
NASA, Boeing, and McDonnell Douglas, who were then working under the
Integrated Wing Design element of NASA's Advanced Subsonic Technology
Program.
The work being planned involved extensive use of CFD, and the
possibility of collaborative analyses by many organizations.
It was thus necessary to establish a common data format suitable to meet
the needs of production CFD tools in the mid- to late-1990s.
This format would be used to enable interchange of data among different
CFD-related tools and different computing platforms, and to provide a
mechanism for archive and retrieval of CFD data

At that time, the \textit{de facto} standard for CFD data was the file
format used by Plot3D, a flow visualization program developed at NASA
Ames.
However, the Plot3D format was developed to expedite post-processing,
and was never intended as a general-purpose standard for the storage
and exchange of CFD data.
By the early 1990s, as CFD became more sophisticated, the limitations of
the Plot3D format as a general-purpose standard had become apparent.
It had no provisions for storing data associated with modern CFD
technology, such as unstructured grids, turbulence models based on
solutions of partial differential equations, and flows with multiple
chemical species.

Individual organizations were overcoming these limitations by defining
extensions to the Plot3D format to meet their needs.
These extensions were not coordinated among different organizations, and
therefore data stored in these extended formats generally could not be
utilized outside the originating organization.
Further, the Plot3D format was not self-documenting, and it was
necessary to rely on file-naming conventions or external notes to
maintain awareness of the flow conditions and analyzed geometry of each
PLOT3D data file.

To overcome these limitations, several database options were considered
by the NASA~/ Boeing~/ McDonnell Douglas team from December 1994 to
March 1995.
In March 1995, the decision was made to build a new data standard called
CGNS (Complex Geometry Navier Stokes).
This standard was a ``clean sheet'' development, but it was heavily
influenced by the McDonnell Douglas Common File Format (CFF), which had
been established and deployed in 1989 and heavily revised in 1992.

Agreement was reached to develop CGNS at Boeing, under NASA Contract
NAS1-20267, with active participation by a team of CFD researchers from

\begin{itemize*}
   \item NASA Langley Research Center
   \item NASA Glenn Research Center
   \item NASA Ames Research Center
   \item McDonnell Douglas Corporation (now Boeing St. Louis)
   \item Boeing Commercial Airplane Group (Seattle) Aerodynamics
   \item Boeing Commercial Airplane Group (Seattle) Propulsion
   \item ICEM CFD Engineering Corporation
\end{itemize*}

The original development work itself occurred from 1995--1998, and
consisted of essentially four activities:

\begin{itemize*}
\item Development of the standalone database manager (the ADF Core)
\item Identification and layout of the CFD data (the SIDS and File
      Mapping)
\item Development of an API for use in applications codes (the
      Mid-Level Library)
\item Demonstration of system capability
\end{itemize*}

\subsubsection{Development of ADF}

As noted earlier, the decision to develop CGNS was made after an
examination of several existing formats.
Because of its hierarchical structure, the CFF (Common File Format) that
was then in use at McDonnell Douglas was judged the most compatible with
the SIDS.
However, CFF had known limitations, and it was decided to adopt the CFF
structure but recast it in a more general form.
This gave the developers control over the software and facilitated the
development of a robust, portable, reliable, and flexible database
manager, which could be distributed freely within the system.

These low-level routines were developed during 1995, and make up the
ADF (Advanced Data Format) Core.
ADF was written and extensively tested at Boeing by a small subteam
familiar with standard software development practices, and it may now be
considered a stable product.

\subsubsection{Identification and layout of the CFD data}

Another task, undertaken in parallel with the development of ADF, was to
identify the data associated with CFD and establish a means of encoding
it within the ADF format.
This task was the primary concern of the national team, with Boeing
proposing the standards and the remainder of the team serving to
critique, test, and improve them.
This began with the development of the SIDS and its language.
The hierarchical nature of the SIDS was used to suggest the SIDS-to-ADF
File Mapping conventions, which were developed slightly after the SIDS
but along the same track.

\subsubsection{Development of an API}

At a review in June 1997, the CGNS team (NASA, Boeing, and McDonnell
Douglas) determined that additional support would be required to produce
an adequate Mid-Level Library.
Subcontracts were issued to the ICEM CFD Engineering Company, in
Berkeley, CA, following this decision.
ICEM CFD in effect became the lead organization for the development of
the Mid-Level Library.

Also at this time, the acronym ``CGNS'' was re-defined to mean ``CFD
Generalized Notation System'', which was more in keeping with the evolved
goals of this project.

Version 1.0 of the Mid-Level Library, which met the original goal of
supporting structured multi-block analysis codes, was released in May
1998.
This set of higher-level routines enables a user to implement the CFD
standards defined by the SIDS, without specific knowledge of the File
Mapping or the ADF Core.
This activity was crucial to the acceptance and implementation of CGNS
among the CFD community.

NASA and the informal CGNS committee determined that there was no
need for export authority so the CGNS standard, the ADF Core and the
Mid-Level Library, and all supporting documentation could be distributed
worldwide as freeware.
Appropriate legal reviews and approvals were obtained at both NASA and
Boeing to validate this decision.

\subsubsection{Demonstration of system capability}

\paragraph{Prototypes.}
To test the CGNS data recording and retrieval mechanisms and to
gain experience with the installation of the software into common
applications codes, the developers modified two CFD solvers (namely,
NPARC and TLNS3D) to operate in the CGNS environment.
In addition, NASA Langley Research Center modified CFL3D similarly.
These modifications were made early in the project and predate the
Mid-Level Library and some later portions of the SIDS.

Separate data cases were prepared for each of the prototypes.
Each prototype proved capable of starting, exiting, and restarting its
data case as expected.
In addition, in all three cases, transfer of the CGNS data between
workstation (SGI) and mainframe (CRAY) platforms was successfully
demonstrated.

Dissimilarity of the \emph{content} of the data required by NPARC
and the other two codes (i.e., nodal vs. cell-centered data) prevented
restarting NPARC from TLNS3D/CFL3D data and vice versa.
However, this type of restart was demonstrated between TLNS3D and CFL3D.

These prototypes were limited in certain ways and, therefore, were not
suitable as a final implementation of their respective capabilities in
CGNS.
The principal limitations are:

\begin{itemize}
\item The prototypes implement only the grid coordinates, flow solution,
      boundary condition, and 1-to-1 interface connectivity portions of
      the CGNS data specification.
\item No attempt was made to modify the internal structure of any of
      the codes in order to improve compatibility with CGNS data
      organization.
\item The prototypes accessed the ADF files primarily using routines
      at the ADF Core level.
      There were some higher-level routines included, and these, in many
      cases, suggested the content of the CGNS Mid-Level Libraries.
      But the high-level prototype routines often intermingled
      ADF functions with CGNS functions and, in some cases, were
      code-specific or dependent on internal directives to make them so.
      They were thus less broadly applicable than the current Mid-Level
      Library routines.
\item These codes exercised only a portion of the CGNS boundary
      condition specifications. 
\item The prototypes sometimes incorporated extra nodes into their ADF
      files to carry code-specific data.
      This practice arose partially from the lack of complete CGNS data
      specifications at the time the prototypes were written, but it
      resulted also, in part, because the current code input structure
      required a different (but equivalent) form of the CGNS data, and
      the developers opted to duplicate it within the CGNS database.
\end{itemize}

The general lesson learned from the construction of the prototypes
was that professional programmers had no conceptual difficulty in
implementing CGNS at the ADF Core level.
But the resulting code was cumbersome, and the development of the
Mid-Level Library was needed to facilitate dissemination of CGNS among
those disinclined to work with the ADF Core.

\paragraph{System Demonstrator.}
As the CGNS effort progressed, an additional activity was undertaken
to demonstrate the capability of the mature system.
Known as the ``system demonstrator'', this tested the ability of CGNS to
transfer data seamlessly between applications that had never operated
together before.
The system demonstrator used most of that portion of the currently
existing SIDS that has been included in the Mid-Level Library.

The geometry chosen for the test was a high-lift configuration known as
the trapezoidal wing.
This is a multielement airfoil with a full-span slat and flap, and a
generic fuselage.

Three separate CGNS-compatible application codes were involved in the
system demonstrator.

\begin{Ventryic}{OVERFLOW}
\item [GMAN]
      A pre-processor and grid generation tool developed by Boeing
      St. Louis
\item [OVERFLOW]
      A CFD flow solver developed by NASA Ames
\item [Visual3]
      A flow visualization tool from ICEM CFD, originally developed at
      MIT
\end{Ventryic}

The system demonstrator consisted of the following tasks:

\begin{enumerate}
\item The grid was generated using NASA Ames grid tools and written
      as a Plot3D file.
\item The grid file was sent to Boeing St. Louis, where it was processed
      by a locally-modified version of GMAN.
      GMAN calculated the grid connectivity information, and wrote the
      grid and connectivity data into a CGNS database.
      Boundary conditions were also added to the CGNS database at this
      time.
\item The CGNS file was returned to NASA Ames, where it was read by the
      newly-modified OVERFLOW code.
      (Some iteration was necessary here, because the definition of
      overset holes used by GMAN differs from that normally expected by
      OVERFLOW.
      The CGNS file was intact, and served, as intended, to highlight the
      discrepancy.)
      OVERFLOW computed the flow field, writing the results into the
      CGNS database.
\item The CGNS file was next sent to ICEM CFD.
      There, the CGNS database was read and displayed using Visual3.
\end{enumerate}

The system demonstrator involved significant cross-platform transfer
among various workstation and mainframe computing facilities.
The results indicated that universal data exchange was well supported by
CGNS.

\subsection{Subsequent Development}

\subsubsection{Software}

Since the release of Version 1.0 of the CGNS software in May 1998,
several improvements and extensions have been made to the SIDS, the File
Mapping, and the Mid-Level Library.

\begin{itemize}
\item Version 1.1, released in June 1999, added support for
      unstructured grids, and for associating CAD geometric entities
      with grid surfaces.
\item Version 2.0, released in December 2000, added support for
      moving and/or deforming grids, and for iterative/time-dependent
      data.
\item Version 2.1, released in May 2002, added support for
      user-defined data arrays, chemistry, and linked nodes.
\item Version 2.2, released in May 2003, added support for axisymmetry
      information, rotating coordinates, special properties associated
      with particular grid connectivity patches, such as periodicity or
      averaging, special properties associated with particular boundary
      condition patches, such as wall functions and bleed, and gravity.
\item Version 2.3, released in Jan 2004, restored the capability to
      specify a boundary condition patch using \fort{ElementRange} or
      \fort{ElementList}.
\item Version 2.4, released in Aug 2005, added support for describing
      electric field, magnetic field, and conductivity models used
      in electromagnetic flows, specification of units for electric
      current, substance amount, and luminous intensity, more flexible
      specification of boundary condition locations, and additional
      user-defined data.
      Also new in Version 2.4 is the ability to choose, at build time,
      either ADF or HDF5 as the underlying database manager.
\item Version 2.5, released in Sep 2007, added Mid-Level Library
      functions to check file validity, configure some internal CGNS
      library options, and provide alternate ways to access a node.
      Some changes were also made in the use of functions for partial
      writes of coordinate, element, and solution data.
      Support was also added for building the CGNS library as a DLL under
      Windows.
\end{itemize}
A more detailed list of revisions to the CGNS software is available at
the CGNS web site.

ICEM CFD served as the focal point for CGNS software development through
the release of Version 2.0, plus the changes in Version 2.5, using
internal company resources.
New features in Version 2.1 and 2.4 were added by Intelligent Light,
with funding from NASA Langley.
Eagle Aeronautics, with additional participation by ICEM CFD, added the
new features in Version 2.2, again with funding from NASA Langley.
Support for HDF as the underlying database in Version 2.4 was added
by personnel from ONERA, ICEM CFD, and the U. S. Air Force Arnold
Engineering Development Center.

\subsubsection{Management and Support}

In 1998 NASA announced that the Advanced Subsonic Technology program,
which had funded the initial CGNS development, would end in September
1999.
Several organizations interested in the continued development of CGNS
met in May 1999 to discuss options.
The decision was made to create the CGNS Steering Committee, a voluntary
public forum made up of international representatives from government
and private industry.

The Steering Committee is responsible for coordinating the further
development and dissemination of the CGNS standard and its supporting
software and documentation.
In Jan 2000, the CGNS Steering Committee became a Sub-committee of the
AIAA CFD Committee on Standards.
Additional details about the mission and responsibilities of the
Steering Committee, and its organization, are in the \textit{CGNS Steering
Committee Charter}.

The basic CGNS documentation has of course been updated to reflect the
software changes described above.
In addition, all the documentation has been converted to LaTeX (used
to create PDF versions for printing), and to HTML (for interactive
use).
A new document, \textit{A User's Guide to CGNS}, was made available in
October 2001, and is very useful for new and prospective users of the
Mid-Level Library.

To encourage communication between CGNS users, a mailing list called
CGNStalk was created in Oct 2000.
Instructions for subscribing are available at the CGNS web site.

\subsection{Current Status}

Since the initial release of the Mid-Level Library in May 1998,
interest in CGNS has continued to grow throughout the global CFD
community.
The system is being used by engineers and scientists in academia,
industry, and government.
By January 2002, 591 users from more than 25 countries had formally
registered at the CGNS web site.
As of September 2007, the CGNStalk mailing list had 224 participants
from 20 different countries and at least 79 different organizations.

\subsubsection{Software}

The current ``production'' release of CGNS is Version 2.5, released in
September 2007.

Several extensions to CGNS have been formally proposed.
Documentation supporting these extensions, and information on their
current status, is available via the ``Proposals for Extensions'' link at
the CGNS web site.

\subsubsection{Standards}

Between 1999 and 2002, an effort was spearheaded by Boeing to establish
an ISO-STEP standard for the representation, storage, and exchange of
digital data in fluid dynamics based on the CGNS standard.
Unfortunately, the effort had to be curtailed because of budget
problems.
It was subsequently decided that an existing ISO standard on finite
element solid mechanics would be rewritten and submitted to include CGNS
as well as an integrated engineering analysis framework.

As part of its role as a sub-committee of the AIAA CFD Committee
on Standards, the CGNS Steering Committee is also involved in the
development of the AIAA Recommended Practice for the storage of CFD
data.
The Recommended Practice consists of the CGNS Standard Interface
Data Structures (SIDS) document, reformatted to conform to AIAA's
requirements.
The current AIAA Recommended Practice (corresponding to CGNS Version
2.4) is available at the
\href{http://www.aiaa.org/content.cfm?pageid=363&id=1657}{AIAA
Online Store}, and as a
\href{http://www.grc.nasa.gov/www/cgns/sids/aiaa/R_101A_2005.pdf}{PDF file}
(1.01M, 200 pages) at the CGNS Documentation web site.

\subsection{Noteworthy Contributions}

The entire CGNS team deserves credit for its accomplishments.
The project might easily be cited as a prime example of effective
voluntary cooperation between government and industry.
Every member has contributed ideas, enthusiasm, and improvements, and
enabled us to represent a cross-section of knowledge and practice from
the CFD community.

It would not be practical to outline each member's contributions to a
project of this complexity.
Nevertheless, it is appropriate to note the contributions of a few
individuals who dedicated special effort during the crucial stages
leading to the initial release of CGNS in May 1998.

\begin{itemize}
\item Steve Allmaras conceived, pursued, and completed the SIDS and
      developed the language in which it is written.
\item Tom Dickens, Matt Smith, and Wayne Jones designed the ADF Core.
\item Tom Dickens wrote the ADF Core.
\item Dan Owen maintained and tuned the ADF Core after Tom left the group.
\item Chuck Keagle designed and executed stringent testing procedures
      for the ADF Core; that is to say, he did everything he could think
      of to break it.
\item Diane Poirier designed and wrote the CGNS Midlevel Library and
      drafted the original proposals for unstructured grid.
      She and Alan Magnuson drafted the proposals for CAD-to-geometry
      specifications.
\item Gary Shurtleff wrote the TLNS3D and NPARC prototypes and provided
      us with what was, for a long time, the only examples of working
      CGNS software.
\item Chris Rumsey served as liaison with NASA Langley, responding in
      detail to requests for criticism and improvements and testing the
      system by writing the CFL3D prototype.
\item Cetin Kiris converted the OVERFLOW code to CGNS, twice.
\item Wayne Jones tested our abstractions with real code, which found
      its way into the ADF Midlevel Library and the CGNS Toolkit.
\item Matt Smith wrote much of the File Mapping and ADF Core documents
      and, with knowledge of both CFD and software design, brought good
      sense to bear on proposals that needed it.
\item Ray Cosner's retrospective vantage point made him a reliable
      supporter who was often able to move things forward when progress
      slowed.
\item Susan Jacob supplied initial guidance as a Productivity+ coach and
      got us all moving in (more or less) the same direction.
\item Doug McCarthy led the project to completion.
\item Shay Gould made the documentation intelligible and presentable.
\end{itemize}

\noindent
And, a particularly special mention:

\begin{itemize}
\item Ben Paul secured the initial funding and shepherded the project
      through the early phases of the contract.
      We will never forget the many meetings Ben held to make sure that
      he thought that he knew that we thought that we knew what we needed
      to do, and that we were doing it.
      His patience and good-natured perseverance could not have been
      replaced.
\end{itemize}


\end{document}
